\centerline{\bf \Huge Вторая МатДрака}

\textit{ВСЕ участники команды, занявшей I место, получают по 4 балла в рейтинговую таблицу и право на 4 попытки ввода кода.}

\textit{ВСЕ участники команды, занявшей II место, получают по 3 балла в рейтинговую таблицу и право на 3 попытки ввода кода.}

\textit{ВСЕ участники команды, занявшей III место, получают по 2 балла в рейтинговую таблицу и право на 2 попытки ввода кода.}

\textit{ВСЕ участники команды, занявшей IV место, получают по 1 баллу в рейтинговую таблицу и право на 1 попытку ввода кода.}

\textbf{Ученик, решивший задачу, получает 2 балла в рейтинговую таблицу и право на 2 попытки ввода кода; если решали двое, то по 1 баллу и право на 1 попытку ввода кода.}

\problem{\textbf{1}} Найдите все решения ребуса: А + А + Б = 24 - 3$\times$Б \\

\problem{\textbf{2}} Найдите все решения ребуса: ГОЛ = 2026 - ГОЛЫ\\

\problem{\textbf{3}} Найдите все решения ребуса: НИКТО - КТО + ЧТО = 12345\\

\problem{\textbf{4}} Найдите все решения ребуса: ДА + НЕ = 67\\

\problem{\textbf{5}} Найдите все решения ребуса:  АРКА + РКА + КА + А = 2014.\\

\problem{\textbf{6}} Найдите все решения ребуса:  ЛЕТО + ЛЕС = 2011.\\


\problem{\textbf{7}} У каждого из тридцати шестиклассников есть одна ручка, один карандаш и одна линейка. После их участия в олимпиаде оказалось, что 26 учеников потеряли ручку, 23 – линейку и 21 – карандаш. Найдите наименьшее возможное количество шестиклассников, потерявших все три предмета.

\problem{\textbf{8}} Из 100 студентов университета английский язык знают 28 студентов, немецкий — 30, французский — 42, английский и немецкий — 8, английский и французский — 10, немецкий и французский — 5, все три языка знают 3 студента. Сколько студентов не знают ни одного из трех языков?\\

\problem{\textbf{9}} Каждый ученик класса - может быть девочкой, блондином или математиком. В классе 20 девочек, из них 12
блондинок, но одна блондинка математик. Всего в классе 24 ученика - блондина, среди них математиков 12, а всего
математиков (мальчиков и девочек) 17, из них 6 девочек. Сколько учеников в данном классе?\\

\problem{\textbf{10}} Составьте из цифр 1, 2, 3, 4, 5, 6, 7, 8, 9 магический квадрат, то есть разместите их в таблице $3\times3$ так, чтобы суммы чисел по строкам, столбцам и двум диагоналям были одинаковы.\\

\problem{\textbf{11}} Можно ли квадрат разрезать на 9 квадратов и раскрасить их так, чтобы получились 1 белый, 3 серых и 5 чёрных квадратов, причём одноцветные квадраты были бы равны, а разноцветные квадраты – не равны?\\

\problem{\textbf{12}} Можно ли поставить в ряд все натуральные числа от 1 до 100 так, чтобы каждые два соседних числа отличались либо на 2, либо в два раза?\\

\problem{\textbf{13}} Использовав каждую из цифр от 0 до 9 ровно по разу, запишите 5 ненулевых чисел так, чтобы каждое делилось на предыдущее.\\

\problem{\textbf{14}} В семизначном числе первую и последние две цифры заменили
звёздочками: *2026**. Известно, что число делится на 24. Восстановите число\\


\problem{\textbf{15}} Найдите количество четырёхзначных чисел, у которых третья цифра
меньше четвёртой на 2.



